% F2 -- Anatomy of a cell name. A schema: almost no text inside the drawing.
\begin{figure}[!htbp]
\centering
\fitwidth{%
\begin{tikzpicture}[font=\small,
  seg/.style={draw=clrule,minimum height=1.0cm,inner xsep=8pt,align=center},
  tag/.style={align=center,font=\footnotesize,text width=3.0cm},
  ar/.style={-{Stealth[length=4pt]},clrule},
  chip/.style={draw=clrule,fill=white,rounded corners=2pt,inner sep=4pt,
               font=\footnotesize\ttfamily},
]
% ---------------- the name, taken apart ----------------
\node[seg,fill=clmodel!14] (t)                  at (2.2,0) {\ttfamily\LARGE B};
\node[seg,fill=clink!18,right=0pt of t]  (root) {\ttfamily\LARGE SOCHO};
\node[seg,fill=clmove!18,right=0pt of root] (z) {\ttfamily\LARGE 1};

\node[tag,anchor=north] (nt)    at (0.9,-1.05)
  {\textbf{2G or 3G}\\ \cid{B}, \cid{D} = 2G\\ \cid{U}, \cid{V} = 3G};
\node[tag,anchor=north] (nroot) at (4.5,-1.05)
  {\textbf{which mast}\\ \emph{this, and only this,}\\ \emph{is a place}};
\node[tag,anchor=north] (nz)    at (8.1,-1.05)
  {\textbf{which antenna}\\ of that mast\\ (one digit, or three)};

\draw[ar] (t.south)    -- ++(0,-0.25) -| (nt.north);
\draw[ar] (root.south) -- ++(0,-0.25) -| (nroot.north);
\draw[ar] (z.south)    -- ++(0,-0.25) -| (nz.north);

% ---------------- four names, one mast ----------------
\begin{scope}[xshift=11.35cm,yshift=0.65cm]
  \node[chip] (c1) at (0,0)       {BKVCHE1};
  \node[chip] (c2) at (2.15,0)    {BKVCHE4};
  \node[chip] (c3) at (0,-0.72)   {UKVCHE102};
  \node[chip] (c4) at (2.15,-0.72){UKVCHE201};
  \node[draw=clink,fill=clink!12,rounded corners=2pt,inner sep=5pt,
        font=\footnotesize,text width=4.4cm,align=center] (share) at (1.07,-2.05)
        {four names, one mast: \cid{KVCHE}};
  \draw[ar,clink] (c3.south) -- ++(0,-0.30) -| (share.north);
  \draw[ar,clink] (c4.south) -- ++(0,-0.30) -| (share.north);
\end{scope}
\end{tikzpicture}}
\caption{What the name of an antenna contains}
\label{fig:cellid-anatomy}
\figtext{%
\figpara{The name is not a code, it is a description.}{On the left, the antenna
name \cid{BSOCHO1} taken apart. The first letter says which generation of radio
served the phone, 2G or 3G. The middle letters name the mast, and that middle
part is the only piece of the name that corresponds to a location on a map. The
last digits say which antenna of that mast the phone was on, and since a mast
carries three or four antennas pointing in different directions, those digits
change without anybody going anywhere.}

\figpara{The consequence, on the right.}{The same mast in the town of Cheb
appears in the data under four different names, because it carries two
generations of radio and several antennas each. A phone switching from
\cid{BKVCHE1} to \cid{UKVCHE102} has changed name and not changed place.}

\figpara{Why this decides how everything else is measured.}{On the day used for
the analysis of Chapter~\ref{ch:groups}, \textbf{12.4 out of every 100}
back-to-back records are exactly that: a new antenna name, the same mast, and a
person who did not move. Counting those as movement would inflate every mobility
figure in this report by about a quarter and would scramble the four kinds of
user of Chapter~\ref{ch:groups}, whose entire purpose is to separate people who
move from people who do not. So whenever this report measures \emph{movement},
it compares masts. What the model has to \emph{predict}, on the other hand, stays
the full antenna name, because that is the question the operator asked.}}
\end{figure}
